% ID: SEL-D0009-BC-E05
% ESTADO: validada. S010; tres pasadas del mismo autor y comprobación matemática; sin segundo revisor.
\textbf{Datos y descomposición en sucesos disjuntos.}\par\nopagebreak[4]
$P(A)=1-P(A^c)=0.65$. Como $A$ se descompone en $A\setminus B$ y $A\cap B$,
\[
P(A\cap B)=0.65-0.30=0.35.
\]
Por tanto $P(B\setminus A)=0.55-0.35=0.20$.

\textbf{a) Unión.}\par\nopagebreak[4]
La inclusión-exclusión evita contar dos veces la intersección:
\[
\boxed{P(A\cup B)=0.65+0.55-0.35=0.85.}
\]
\textbf{b) Probabilidad condicionada.}\par\nopagebreak[4]
Se restringe el espacio al suceso $A^c$, que tiene probabilidad positiva:
\[
\boxed{P(B\mid A^c)=\frac{P(B\cap A^c)}{P(A^c)}=\frac{0.20}{0.35}=\frac47.}
\]
\textbf{c) Ninguno de los dos.}\par\nopagebreak[4]
Por las leyes de De Morgan,
\[
\boxed{P(A^c\cap B^c)=1-P(A\cup B)=0.15.}
\]
\textbf{d) Independencia.}\par\nopagebreak[4]
Serían independientes si $P(A\cap B)=P(A)P(B)$. Sin redondear:
\[
0.35\ne0.65\cdot0.55=0.3575.
\]
\textbf{No son independientes}. Como comprobación, las cuatro regiones disjuntas tienen probabilidades $0.35$, $0.30$, $0.20$ y $0.15$: son no negativas y suman $1$.
